\chapter{Análisis de requisitos}\label{requisitos}

Este apartado estudiará los requisitos del sistema, tanto los funcionales como los no funcionales, lo que delimitará el alcance del proyecto con mayor grado de detalle.

\section{Requisitos funcionales}

Describen cómo debe comportarse la aplicación cara al usuario.

\begin{enumerate}
    \item Como usuario de la aplicación, quiero poder visitar la implementación desde la web.

    \item Como usuario, quiero poder ver una explicación de cada parte del protocolo (prueba de identidad cero-conocimiento, intercambio de claves y encriptación/desencriptación).

    \item Como usuario, quiero que estas explicaciones se realicen a través de unas presentaciones que se reproduzcan automáticamente, pero que esta reproducción automática se pause cuando interactúe con ellas.
\end{enumerate}

\section{Requisitos no funcionales}

Describen otras características del proyecto que no implican cambios en el funcionamiento de este.

\subsection{Accesibilidad}

\begin{enumerate}
    \item Como administrador, quiero que la aplicación cumpla los requisitos mínimos de accesibilidad.

    \item Como administrador, quiero el sitio esté disponible tanto en inglés como en español, además de ser fácilmente escalable a la hora de añadir más idiomas.
\end{enumerate}

\subsection{Presentación}

\begin{enumerate}
    \item Como usuario, quiero que el sitio sea responsivo de manera que no tenga demasiados problemas para consultar el sitio sin importar el dispositivo desde donde entre.

    \item Como administrador, quiero que la visualización de los protocolos sólo sea accesible desde ordenadores y/o dispositivos que tengan un mínimo de dimensión, ya que un móvil o similar distorsiona demasiado las proporciones para que se vean con claridad las explicaciones.
\end{enumerate}